% MT026 / explanation — canonical editable content.
横轴\(z=yf(x)\)是有符号分数，\(z<0\)分错、\(z>0\)分对。本图最低点在1，两侧是直线，形成V形；图中左右斜率可以不同，不能未经题设就把它们设成相同数值。


当\(z<1\)，把分数往右移动会降低损失，所以模型会纠正错分点，也会推动间隔不足的正确点。由于这侧是直线，错分从很远变成更远时，斜率大小不继续增长；比起平方形损失，它不会仅因负间隔的绝对值更大，就给出越来越强的更新信号。


当\(z>1\)，曲线又往上走。\(z=1\)时损失最低，向右越走损失越大；虽然这些点已经明确分对，模型仍会把它们往1拉。它追求的是分数接近固定目标1，而非只要求分类正确且间隔足够。


所以要分别评价两侧：左边线性增长缓和了极端错分影响，右边却继续惩罚容易的正确样本。原解认为后一种压力对间隔分类不理想。与\StudyRef{Answers}{mt008}{MT008}相比，两题都偏好\(z=1\)，但本题两侧斜率有界，不能照搬“平方增长导致极端梯度”的解释。实际准确率仍需根据数据与模型验证。


\textbf{出处：} 2021A期中 Q13，原卷第4页（10分）


\textbf{答案依据：} 2021A期中 Q13 随卷／配套参考答案；条件与解释经复核。
